% Pista ana-24j-q5 · Anàlisi
% Procedència: PAU Catalunya 2024 · Convocatòria de juny · Sèrie 1 · Qüestió 5
\documentclass[11pt,a4paper]{article}
\usepackage{pista}
\pistainfo{Catalunya 2024}{Convocatòria de juny}{Sèrie 1}

\begin{document}

\section*{\textcolor{blauPista}{Qüestió 5}}

\noindent Volem construir un petit cobert de fusta de $6\,\text{m}^{3}$ de volum, en forma de prisma rectangular, adossat a la paret lateral d'una casa, per a guardar-hi llenya. Només cal construir, per tant, el sostre i tres parets (la paret del fons del cobert és la de la casa a la qual està adossat). A més, volem que el cobert mesuri el triple d'amplària que de fondària. Cada metre quadrat de paret té un cost de construcció de $30\,\text{\euro}$ i el sostre costa $50\,\text{\euro}$ per metre quadrat. Un cop construït el cobert, afegir-hi una porta té un cost fix de $35\,\text{\euro}$.

\subsection*{\textit{a)} \normalfont Comproveu que el cost de construcció del cobert ve donat per la funció
\[
C(x) = \dfrac{300}{x} + 150x^{2} + 35,
\]
on $x$ és la fondària del cobert en metres. \hspace*{\fill}\textnormal{[1,25~punts]}}

\pista{En primer lloc, anomena $h$ l'altura del cobert. L'amplària és $3x$, i el volum d'un prisma es calcula seguint aquesta fórmula:
\[
V = a \cdot b \cdot h,
\]
on $a$ i $b$ són els costats de la base i $h$ és l'altura. Imposa que el volum valgui $6$ i aïlla $h$.

\medskip
A continuació, calcula el cost de cada superfície que cal construir. Per exemple, la paret de davant fa $3x$ d'amplària i $h$ d'altura, i per tant costa $30 \cdot 3x \cdot h$ euros. Fes el mateix amb les dues parets dels costats i amb el sostre, que costa $50\,\text{\euro}$ per metre quadrat.

\medskip
Finalment, suma-ho tot, amb els $35\,\text{\euro}$ de la porta, i substitueix-hi $h$. Arribaràs a l'expressió $C(x)$ de l'enunciat.}

\subsection*{\textit{b)} \normalfont Calculeu quines han de ser les dimensions del cobert per tal que el cost de construcció sigui mínim i justifiqueu la resposta. Quin és aquest cost? \hspace*{\fill}\textnormal{[1,25~punts]}}

\pista{En primer lloc, deriva $C(x)$. Recorda que si escrius $\frac{300}{x}$ en forma de potència, $300\,x^{-1}$, pots derivar-la fàcilment. A continuació, resol $C'(x) = 0$. Per treure la $x$ del denominador, multiplica els dos costats de l'equació per $x^{2}$. Comprova amb el signe de $C''$ que és un mínim.

\medskip
Finalment, tingues en compte que la teva resposta ha d'incloure, d'una banda, les tres mides del cobert; i d'altra banda, el cost total. Per a l'altura $h$, fes servir l'expressió que has trobat a l'apartat \textit{a)}.}

\end{document}
